\section{Geometric and group-theoretic preliminaries}\label{sec:preliminaries}

We collect the geometric input that detects a forbidden Bogomolov
sheaf and the group lemmas that will convert a finite Albanese fibre
image into virtual abelianity.

All manifolds are connected unless stated otherwise. A fibration is a
dominant meromorphic map with connected general fibre. We resolve its
graph when making differential or fundamental-group arguments. A smooth
compact space in Fujiki's class \(\mathcal C\) is replaced, when needed,
by a compact K\"ahler modification. Smooth compact bimeromorphic models
have the same fundamental group. Fibre groups always mean groups of
smooth models of general fibres; their images are understood up to
base-point transport. The image of a general fibre group is normal in
the total fundamental group: over a smooth fibration locus this follows
from the homotopy exact sequence, and the fundamental group of the open
total space surjects onto that of the total space.

\subsection{Specialness and quotient reductions}

We recall the precise established inputs used in the proof.

\begin{theorem}[Bogomolov--Campana]\label{thm:specialness-inputs}
For smooth compact K\"ahler manifolds, specialness in the definition of
Section~\ref{sec:introduction} is bimeromorphically invariant and is
preserved by connected finite \'etale covers. The Albanese morphism
of a special manifold is surjective with connected fibres. Moreover, if
a line bundle \(L\) has a nonzero morphism to \(\Omega_X^p\), then
\(\kappa(X,L)\le p\).
\end{theorem}

The Bogomolov-sheaf characterization and stability under finite
\'etale covers are given in
\cite[author's version, Theorems~2.26 and~5.12]{Campana2004};
the remaining specialness properties are developed there. The Iitaka
bound originates in Bogomolov's projective work \cite{Bogomolov1979};
for its compact K\"ahler formulation, see \cite{Campana2015Survey}.
The Albanese statement is recorded with its correction in
\cite[Theorem~1.2]{CampanaCorrection}. Related projective
stability results are developed in \cite{CampanaClaudon}. We only use
surjectivity and connectedness from that correction. The meridian
statement needed for arbitrary, possibly nonprojective Albanese tori
will be proved in Lemma~\ref{lem:albanese-exactness}.
Saturating a rank-one subsheaf and taking its double dual does not
decrease its Iitaka dimension. Thus the line-bundle definition agrees
with the saturated Bogomolov-sheaf formulation.

\begin{lemma}[Effective divisors on a torus]\label{lem:torus-divisor}
Let \(D\ne0\) be an effective rational divisor on a compact complex
torus \(A\). There are a quotient homomorphism \(p:A\to B\) with
connected fibres, where \(B\) is a positive-dimensional abelian
variety, and an ample effective rational divisor \(D_B\) on \(B\)
such that \(D=p^*D_B\). In particular, \(\kappa(A,D)>0\).
\end{lemma}
\begin{proof}
We give the standard torus argument; see also \cite{Debarre1999}.
After multiplying \(D\) by a positive integer, assume it is integral.
Average its positive integration current \([D]\) over translations of
\(A\). The result is a smooth translation-invariant semipositive
\((1,1)\)-form \(\theta\) representing
\(c_1(\mathcal O_A(D))\). Write \(A=V/\Lambda\). The alternating
form of \(\theta\) is integral on \(\Lambda\), so its radical \(W\)
is defined over \(\mathbb Q\). It is a complex subspace because
\(\theta\) has type \((1,1)\). Thus
\(K=W/(W\cap\Lambda)\) is a subtorus. Let \(p:A\to B=A/K\) be
the quotient.

On every \(K\)-coset, \(\mathcal O_A(D)\) has zero real first Chern
class. If its defining section is not identically zero on that
coset, its zero divisor must be empty: a nonzero effective divisor
on the positive-dimensional K\"ahler torus \(K\) has strictly
positive pairing with the \((\dim K-1)\)-st power of a K\"ahler class.
For \(K=0\) the same conclusion is immediate. Each coset is
therefore either contained in \(\operatorname{Supp}D\) or disjoint
from it. If \(B\) were a point, this argument with \(K=A\) would
make the defining section of \(D\) nowhere zero, contrary to
\(D\ne0\). Hence \(\dim B>0\).

Translations by the connected group \(K\) consequently preserve
each prime component \(D_i\) of \(D\), since a connected group
cannot permute their finite set.
The image \(\overline D_i=p(D_i)\) is a prime divisor on \(B\):
properness gives an analytic image, and invariance gives
\(\dim\overline D_i=\dim D_i-\dim K=\dim B-1\).
Moreover \(p^{-1}(\overline D_i)=D_i\), and smoothness of \(p\)
gives \(p^*\overline D_i=D_i\) with multiplicity one. Keeping the
coefficients of the \(D_i\) therefore defines an effective integral
divisor \(D_B\) with \(D=p^*D_B\).

Since \(p^*D_B=D\), uniqueness of translation-invariant cohomology
representatives identifies \(c_1(\mathcal O_B(D_B))\) with the
positive definite form induced by \(\theta\) on \(B\). The positivity
criterion for line bundles on a complex torus makes \(D_B\) ample
and \(B\) an abelian variety. Dividing \(D_B\) by the initial integer
proves the rational statement. Ratios of
sections of its ample multiples remain independent after pullback,
so \(\kappa(A,D)\ge\dim B>0\).
\end{proof}

\begin{theorem}[Meromorphic \(\Gamma\)-reduction]\label{thm:gamma-reduction}
Let \(M\) be a smooth compact K\"ahler manifold and let
\(\rho:\pi_1(M)\to P\) be a homomorphism. There is a fibration
\(\gamma_\rho:M\dashrightarrow B_\rho\) whose general fibre has finite
fundamental-group image under \(\rho\), and which contracts every
compact irreducible subvariety through a very general point whose
fundamental group has finite image under \(\rho\). The groups of
subvarieties may be computed on their smooth resolutions. The base has
a compact K\"ahler model.
\end{theorem}

The contraction assertion is Campana's \(\Gamma\)-reduction
\cite[Theorem~3.5 and Remark~3.6]{CampanaGamma}; the arbitrary-quotient
formulation is also stated in
\cite[author's version, Theorem~2.1 and Corollary~2.2]{CCE2015}.
The K\"ahler-model assertion uses separately the standard stability
of Fujiki's class \(\mathcal C\) under meromorphic images: the base
belongs to \(\mathcal C\), and a smooth model admits a compact
K\"ahler modification.
We call \(\dim B_\rho\) the \emph{\(\Gamma\)-dimension} for \(\rho\),
and call \(\rho\) \emph{generically large} when this dimension is
\(\dim M\). The target group need not be linear. The qualifier
``very general'' means outside a countable union of proper analytic
subsets. The contraction criterion also applies when the relevant
subvarieties exist through a full-measure set: intersect this set with
the very general locus in the theorem.

\begin{theorem}[Linear general-type criterion]\label{thm:linear-general-type}
Suppose a compact K\"ahler manifold has a generically large complex
linear representation with torsion-free image and connected simple
semisimple Zariski closure. Then it is of general type.
\end{theorem}

\begin{proof}
Let \(M\) be the manifold, and regard \(\rho\) as a Zariski-dense
representation into its semisimple closure. By
\cite[author's version, Theorem~1 and Theorem~6.3]{CCE2015}, it has
a Shafarevich morphism
\[
 \operatorname{sh}_\rho:M\longrightarrow\operatorname{Sh}_\rho(M)
\]
with normal projective target of general type. This morphism is a
model of the connected \(\Gamma\)-reduction
\cite[author's version, Definition~2.13]{CCE2015}. Generic largeness
gives \(\dim\operatorname{Sh}_\rho(M)=\dim M\). The connected general fibre
is then a point, so \(\operatorname{sh}_\rho\) is bimeromorphic.
Bimeromorphic invariance of Kodaira dimension proves that \(M\) is
of general type.
\end{proof}

\subsection{Root orbifolds and adjoint positivity}

For a smooth compact K\"ahler manifold \(Y\) and a simple normal
crossing divisor \(\sum D_i\) with integers \(m_i\ge1\), let
\(\cY\) be the analytic root orbifold of orders \(m_i\), and write
\[
 \Delta=\sum_i(1-m_i^{-1})D_i.
\]
Locally, if the divisor coordinates are \(z_i\), its charts have
\(z_i=w_i^{m_i}\) with the corresponding product of cyclic groups.
Tensors and smooth metrics are invariant tensors and smooth metrics
upstairs. Its fundamental group is the fundamental group of the
divisor complement modulo the normal subgroup generated by the
\(m_i\)-th powers of meridians. Orders equal to one add no boundary.
Orbifold covering space theory supplies a connected cover for each
subgroup; the cover need not be a manifold.

These orbifolds are K\"ahler. Add to a sufficiently large pullback
of a coarse K\"ahler form the \(\ii\partial\bar\partial\) of local
squared root norms, using smooth metrics on the divisor bundles and
cutoffs. Their root-direction terms are positive near the divisor and
their negative parts away from it are bounded by the coarse metric.
On a cover of a compact orbifold, finitely many uniformizing charts
give uniform local elliptic estimates, bounded-overlap partitions and
cutoffs. All complete-metric differential operators below are their
closed \(L^2\) extensions. We use
\(\dd\dc=\ii\partial\bar\partial\).

\begin{theorem}[Klt K\"ahler results over tori]\label{thm:wang-torus}
Let \(Z\) be a smooth compact K\"ahler manifold, let \(D\) be an
effective rational divisor such that \((Z,D)\) is klt, and let
\(g:Z\to A\) be a fibration onto a complex torus. For a general
smooth fibre \(F\),
\[
 \kappa(Z,K_Z+D)\ge\kappa(F,K_F+D|_F).
\]
If \(\kappa(Z,K_Z+D)=0\), the Albanese morphism of \(Z\) is
surjective with connected fibres.
\end{theorem}

These are Theorems~A(II) and~C of \cite{JWang2021}; the torus is
not required to be projective. We apply them to smooth pairs with
simple normal crossing boundaries and coefficients strictly less
than one.

For a projective smooth pair, we also use the bigness consequence of
orbifold cotangent positivity in \cite[Theorem~7.11]{CampanaPaun2019}:
a big ordinary line bundle injecting into a tensor power of the
orbifold cotangent, on adapted covers, forces the log canonical
divisor to be big. The projectivity hypothesis will be established
before this result is applied in Section~\ref{sec:hilbert}.

\subsection{Finite kernels and finite covers}

\begin{lemma}\label{lem:finite-by-abelian}
Let \(E\) be a finitely generated group in an exact sequence
\[
 1\longrightarrow F\longrightarrow E\longrightarrow\Z^r
 \longrightarrow1
\]
with \(F\) finite. Then \(E\) is virtually abelian and residually
finite. In particular it has a finite-index subgroup disjoint from
\(F\).
\end{lemma}
\begin{proof}
The centralizer \(E_0\) of \(F\) has finite index. Its commutator
subgroup is contained in the finite central group \(F\cap E_0\).
Choose generators \(x_1,\ldots,x_t\) of \(E_0\) and an integer
\(e>0\) annihilating \(F\). Commutators are central, so
\([x_i^e,x_j^e]=[x_i,x_j]^{e^2}=1\).
The subgroup generated by the \(x_i^e\) is abelian and has finite
index: its image has finite index in the finitely generated
abelianization, and the commutator subgroup is finite.
This abelian subgroup has a torsion-free subgroup of finite index.
Its core in \(E\) is a normal free abelian subgroup of finite index.
Multiples of this core have finite index in \(E\) and separate its
nonzero elements; the finite quotient by the core separates elements
outside it. Hence \(E\) is residually finite. Since \(F\) is finite,
intersecting finitely many finite-index subgroups separates all its
nonidentity elements at once.
\end{proof}

\begin{lemma}\label{lem:realize-finite-index}
Suppose \(G\) has a finitely generated normal subgroup \(N\) with
\(G/N\cong\Z^r\). For every finite-index subgroup \(N_0\le N\),
there is a finite-index subgroup \(G_0\le G\) satisfying
\(G_0\cap N\subseteq N_0\).
\end{lemma}
\begin{proof}
A finitely generated group has only finitely many subgroups of any
given finite index. Intersect all automorphic images of the core of
\(N_0\) in \(N\). The result is a characteristic finite-index
subgroup \(K\le N\) contained in \(N_0\), and hence is normal in
\(G\). The group \(G/K\) is finite-by-\(\Z^r\).
By Lemma~\ref{lem:finite-by-abelian}, it has a finite-index subgroup
disjoint from \(N/K\). Its inverse image is a suitable \(G_0\).
\end{proof}

A group is \emph{FAb} if every finite-index subgroup has finite
abelianization. Quotients of FAb groups and their finite-index
subgroups are FAb. A finitely generated complex linear FAb group
with virtually solvable image has finite image: after passage to
finite index its Zariski closure is connected and solvable and hence
triangularizable. The image is then solvable, and applying finite
abelianization successively to the finite-index derived subgroups
shows that the image is finite. This observation is used only for
linear images; it imposes no linearity hypothesis on the original
fundamental group.
